% define your own macros here
\newcommand{\tlam}[2]{\lambda#1.\,#2}
\newcommand{\tLam}[2]{\Lambda#1.\,#2}
\newcommand{\hastype}[3]{#1 \vdash #2 : #3}
\newcommand{\vrel}[1]{\llbracket #1 \rrbracket_{\mathcal{V}}}
\newcommand{\erel}[1]{\llbracket #1 \rrbracket_{\mathcal{E}}}
\newcommand{\clk}{\mathsf{clk}}
\newcommand{\tmonad}[2]{\mathbb{M}\,{#1}\, #2}
\newcommand{\tick}[1]{\uparrow^{#1}}
\newcommand{\tstr}[1]{\mathsf{Stream}\,#1}
\newcommand{\ttree}[1]{\mathsf{Tree^g}\,#1}
\newcommand{\tint}{\mathsf{Int}}
\newcommand{\head}{\mathsf{hd}\,}
\newcommand{\tail}{\mathsf{tl}\,}
\newcommand{\troot}{\mathsf{root^g}\,}
\newcommand{\tleft}{\mathsf{left^g}\,}
\newcommand{\tright}{\mathsf{right^g}\,}
\newcommand{\fix}{\mathsf{fix}\,}
\newcommand{\delay}{\mathsf{next}\,}
\newcommand{\fold}{\mathsf{fold}\,}
\newcommand{\unfold}{\mathsf{unfold}\,}
\newcommand{\ret}{\mathsf{ret}\,}
\newcommand{\case}[5]{\mathsf{case}\,#1\,\mathsf{of}\;
  \mathsf{inl}\,#2\mapsto #3\mid\mathsf{inr}\,#4\mapsto #5}
\newcommand{\store}{\mathsf{store}\,}
\newcommand{\bind}[3]{\mathsf{bind}\, #1= #2 \, \mathsf{in}\, #3}
\newcommand{\release}[3]{\mathsf{release}\, #1 = #2 \, \mathsf{in}\, #3}
\newcommand{\pull}{\mathsf{prepay}\,}
\newcommand{\lete}[3]{\text{let }#1 = #2 \text{ in } #3}
\newcommand{\tuple}[2]{\langle #1, #2\rangle}
\newcommand{\ttuple}[2]{\langle\!\langle #1, #2\rangle\!\rangle}
\newcommand{\lmatch}[5]{\mathsf{match}\,#1\,\mathsf{with}
  \mid nil\mapsto #2\mid#3::#4\mapsto #5}
\newcommand{\isval}[2]{#1\;\mathsf{val}_{#2}}
\newcommand{\colist}[1]{\mathsf{Colist}\, #1}

\newcommand{\Eff} {\ensuremath{\mathit{eff}}}  % example command without arguments
\newcommand{\Pre} {\ensuremath{\mathit{pre}}}  % (again)

% Note that you can easily specify arguments:
% \newcommand{\someMacro}[2] {Argument 1: #1, Argument 2: #2} % example command with two arguments
% you use it via \someMacro{Hello}{World!}


% the following commands are being provided by the amsthm package
% the first parameter states the new environmet's name that can be
% used (due to this definition here) and the second the name that
% will appear in the PDF document
\theoremstyle{definition}
\newtheorem{definition}{Definition}   % well, a formal definition!
\theoremstyle{plain}
\newtheorem{prop}{Proposition} % like a theorem, but less important or evolved
\newtheorem{lem}{Lemma}        % used within a proof of a theorem
\newtheorem{thm}{Theorem}      % well, a theorem! :) important and evolved
\newtheorem{cor}{Corollary}    % basically either a proposition or theorem,
                               %  but one that follows from another theorem.
% There's a lot you can configure about the appearance. If interested,
% open the manual of amsthm or google for tutorials etc. on that package
\newtheorem{appalem}{Lemma}[chapter]
\newtheorem{appaprop}{Proposition}[chapter]
\newtheorem{appathm}{Theorem}[chapter]

\newtheorem{appblem}{Lemma}[chapter]
\newtheorem{appbthm}{Theorem}[chapter]
\newtheorem{appbdef}{Definition}[chapter]


% the following add a symbol to the definition environment to make it more
% clear when a definition ends (as there is no difference in fonts!). From:
% https://tex.stackexchange.com/questions/226334/change-a-amsthm-theorem-ending
\newcommand{\xqed}[1]{%
    \leavevmode\unskip\penalty9999 \hbox{}\nobreak\hfill
    \quad\hbox{\ensuremath{#1}}}
\newcommand{\Endofdef}{\xqed{\blacksquare}}
\newenvironment{defn}[1]{%
    \begin{definition}#1}{%
    \Endofdef\end{definition}%
}
